\section{Introduction}
\subsection{Contexte}
Dans le cadre de notre formation, nous avons développé un système embarqué utilisant un microcontrôleur STM32F411RE et une RaspberryPi 3b+.

\subsection{Objectif}
L'objectif principal était de développer un système fonctionne incluant une STM32 et une RaspberryPi.
Le système est capable de :
\begin{itemize}
    \item Retourner la distance mesurée par le capteur ultrasonique HC-SR04 sur une interface de commande sur la Raspberry Pi.
    \item Controler le servomoteur connecté sur la STM32 depuis une interface de commande sur la Raspberry Pi.
    \item De controler en automatique les deux servomoteurs en fonction de la distance mesurée par le capteur ultrasonique au travers d'un calcul réalisé sur la RaspberryPi.
\end{itemize}

Pour cela, la STM32 doit être configurée pour :
\begin{itemize}
    \item Mesurer une distance avec le capteur HC-SR04 toutes les 250 ms
    \item Communiquer les informations de distance via la liaison série toutes les 250 ms
    \item Recevoir une consigne de position pour le servomoteur via la liaison série et l'appliquer au servomoteur
\end{itemize}

Et la Raspberry Pi doit être configurée pour :
\begin{itemize}
    \item Recevoir la distance mesurée par la STM32 via la liaison série
    \item Envoyer une consigne de position pour le servomoteur à la STM32 via la liaison série
    \item Calculer la position de son servomoteur en fonction de la distance mesurée par le capteur ultrasonique
    \item Envoyer la position calculée au servomoteur de la STM32 via la liaison série
\end{itemize}

